\FifthArticle{14}
The preceding rules can now easily be recast as rules for the use of auxiliary logarithms, since these $A=\log a$.
$B=\log b$, $C=\log c$, can be regarded as the logarithms of the squares of the tangents, cosecants and secants of angles increasing from $45^0$ to $90^0$, or, equivalently, as the logarithms of the squares of the cotangents,
secants and cosecants of angles decreasing from $45^0$ to $0$. Thus
\[
a=\mathrm{tang}\theta^2,\quad \frac{1}{b}=\sin\theta^2,\quad
\frac{1}{c}=\cos\theta^2
\]
for values of $\theta$ in the second octant, or
\[
\frac{1}{a}=\mathrm{tang}\theta^2,\quad
\frac{1}{c}=\sin\theta^2,\quad
\frac{1}{b}=\cos\theta^2
\]
for values of $\theta$ in the first octant.

I combine the complete rules in the following scheme, where, as above,
\[
\lambda=\frac{f^n}{e^{m+n}}
\]
is put.

\Needspace{6\baselineskip}\begin{center}\textit{First Form.}\end{center}
\[
x^{m+n}+ex^m-f=0
\]
\textit{First Case.} $\lambda>2^n$
\[
\lambda=a^{m+n}b^n=a^m c^n=\frac{c^{m+n}}{b^m}
\]
\[
x^{m+n}=\frac{f}{b},\quad x^m=\frac{f}{ec},\quad x^n=ea
\]
\textit{Second Case.} $\lambda<2^n$
\[
\lambda=\frac{b^n}{a^m}=\frac{c^n}{a^{m+n}}=\frac{b^{m+n}}{c^m}
\]
\[
x^{m+n}=\frac{f}{c},\quad x^m=\frac{f}{eb},\quad x^n=\frac{e}{a}
\]

\Needspace{6\baselineskip}\begin{center}\textit{Second Form.}\end{center}
\[
x^{m+n}-ex^m-f=0
\]
\textit{First Case.} $\lambda>2^m$
\SourceLeaf{098}
\[
\lambda=a^{m+n}b^m=a^nc^m=\frac{c^{m+n}}{b^n}
\]
\[
x^{m+n}=fb,\quad x^m=\frac{f}{ea},\quad x^n=ec
\]
\textit{Second Case.}\quad $\lambda<2^m$
\[
\lambda=\frac{b^m}{a^n}=\frac{c^m}{a^{m+n}}=\frac{b^{m+n}}{c^n}
\]
\[
x^{m+n}=fc,\quad x^m=\frac{fa}{e},\quad x^n=eb
\]
\Needspace{6\baselineskip}\begin{center}\textit{Third Form.}\end{center}
\[
x^{m+n}-ex^m+f=0
\]
\textit{First Case.}\quad $\frac{1}{\lambda}<\frac{(m+n)^{m+n}}{m^mn^n}$

No solution at all.

\textit{Second Case.}\quad $\frac{1}{\lambda}=\frac{(m+n)^{m+n}}{m^mn^n}$

Two equal roots, for whose determination one of the equations
\[
x^{m+n}=\frac{fm}{n},\quad x^m=\frac{f(m+n)}{en},\quad x^n=\frac{em}{m+n}
\]
is used.

\textit{Third Case.}\quad $\frac{1}{\lambda}$ greater than $\frac{(m+n)^{m+n}}{m^mn^n}$ but not greater than $2^{m+n}$, and at the same time
$m$ greater than $n$.

Two roots, for which
\[
\frac{1}{\lambda}=a^nb^{m+n}=\frac{c^{m+n}}{a^m}=b^mc^n
\]
\[
x^{m+n}=fa,\quad x^m=\frac{fc}{e},\quad x^n=\frac{e}{b}
\]
\textit{Fourth Case.}\quad The same bounds for $\frac{1}{\lambda}$ as in the third case, but $m$ smaller
than $n$.

Two roots, for which
\[
\frac{1}{\lambda}=a^mb^{m+n}=\frac{c^{m+n}}{a^n}=b^nc^m
\]
\[
x^{m+n}=\frac{f}{a},\quad x^m=\frac{fb}{e},\quad x^n=\frac{e}{c}
\]
\textit{Fifth Case.}\quad $\frac{1}{\lambda}$ greater than $2^{m+n}$.

Two roots, one determined by the formulas of the third case and
the other by those of the fourth.
\SourceLeaf{099}
It may also be noted that in the third case the value of $a$ corresponding to one root is smaller than $\frac{m}{n}$, while that corresponding to the other is greater than $\frac{m}{n}$; in the fourth case the two values of
$a$ stand in a similar relation to $\frac{n}{m}$.

